% 第 2 章 Working With DC Explorer
\bichapter{使用 DC Explorer}{Working With DC Explorer}
\label{chap:working}

DC Explorer 为综合与时序分析提供两种界面：\cmd{de\_shell} 命令行界面（shell）以及 Design Vision 图形用户界面（GUI）。
\cmd{de\_shell} 是纯文本环境，您在命令行提示符下输入命令。
Design Vision 是 Synopsys 综合环境的 GUI，用于可视化设计数据并分析结果。

要了解在 DC Explorer 环境中的标准操作任务，请参阅：
\begin{itemize}
  \item 运行 DC Explorer
  \item 许可证（Licensing）
  \item 使用多核技术
  \item 并行运行命令
  \item 在后台运行命令以改善运行时间
\end{itemize}

% ============================================================
\section{运行 DC Explorer}
\subsection*{Running DC Explorer}

下列各节介绍如何使用 \cmd{de\_shell} 与 Design Vision GUI 运行 DC Explorer：
\begin{itemize}
  \item 启动 DC Explorer
  \item 输入 \cmd{de\_shell} 命令
  \item 中断或终止命令处理
  \item 在 \cmd{de\_shell} 中打开或关闭 GUI
  \item 在命令行获取帮助
  \item Setup 文件
  \item 在日志文件中查找会话信息
  \item 运行 Tcl 脚本
  \item 退出 DC Explorer
\end{itemize}

\subsection{启动 DC Explorer}
\subsubsection*{Starting DC Explorer}

DC Explorer 在 Linux 的 X Window 环境下运行。启动前请确保 \texttt{\$PATH} 变量包含指向 \texttt{bin} 目录的路径。
在 Linux shell 中输入 \cmd{de\_shell} 即可启动 DC Explorer。例如：
\begin{lstlisting}
% de_shell
\end{lstlisting}
默认情况下，该命令以命令行界面（\cmd{de\_shell}）启动工具。

启动时还可在命令行加入其他选项。例如，可以使用：
\begin{itemize}
  \item \opt{-checkout}：检出程序默认之外的许可功能
  \item \opt{-wait}：为检出额外许可证设置等待时间上限
  \item \opt{-f}：在显示初始 \cmd{de\_shell} 提示符之前执行脚本文件
  \item \opt{-x}：在启动时执行一条 \cmd{de\_shell} 语句
  \item \opt{-no\_init}：指定 \cmd{de\_shell} 不执行任何 \texttt{.synopsys\_dc.setup} 启动文件\\
        仅在需要用命令日志或其他脚本文件复现先前 \cmd{de\_shell} 会话时使用此选项。
  \item \opt{-no\_home\_init}：指定不执行 home 目录下的 \texttt{.synopsys\_dc.setup}
  \item \opt{-no\_local\_init}：指定不执行本地目录下的 \texttt{.synopsys\_dc.setup}
  \item \opt{-gui}：以 GUI 启动 \cmd{de\_shell}
  \item \opt{-no\_gui}：不以 GUI 启动 \cmd{de\_shell}
\end{itemize}

详细选项列表见 \cmd{de\_shell} man 页。

启动时，\cmd{de\_shell} 执行下列任务：
\begin{enumerate}
  \item 创建命令日志文件。
  \item 读取并执行 \texttt{.synopsys\_dc.setup} 文件。
  \item 分别执行命令行上由 \opt{-f} 或 \opt{-x} 指定的脚本文件或命令。
  \item 在调用 \cmd{de\_shell} 的窗口中显示程序头与 \cmd{de\_shell} 提示符。程序头列出您站点已获许可的全部功能。
\end{enumerate}

\begin{noteBox}
若以路径形式输入 \cmd{de\_shell} 命令，必须指定完整绝对路径，而不能使用相对路径。
若使用相对路径（如 \texttt{../}），DC Explorer 将无法访问位于根目录中的库。

下列示例给出正确指示包含 DC Explorer 安装的 Synopsys 根目录的方式：
\begin{lstlisting}
/tools/synopsys/2012.06/bin/de_shell
\end{lstlisting}
\end{noteBox}

\subsection{输入 de\_shell 命令}
\subsubsection*{Entering de\_shell Commands}

您通过 \cmd{de\_shell} 命令与 DC Explorer 交互；这些命令基于工具命令语言（Tcl），并包含实现特定 DC Explorer 功能所需的命令扩展。
DC Explorer 命令语言提供与 Linux 命令 shell 类似的能力，包括变量、条件执行与控制流命令。您可以：
\begin{itemize}
  \item 在 \texttt{de\_shell>} 提示符下交互输入单条命令
  \item 在 Design Vision GUI 的控制台命令行上交互输入单条命令
  \item 运行一个或多个 Tcl 命令脚本（包含 \cmd{de\_shell} 命令的文本文件）
\end{itemize}

在 \cmd{de\_shell} 中输入命令的方式与标准 Linux shell 相同。
GUI 打开时，可在控制台命令行输入命令，也可通过菜单界面使用 DC Explorer 命令。

在 \cmd{de\_shell} 中输入命令：
\begin{enumerate}
  \item 在命令行上键入命令。
  \item 按 Return。
\end{enumerate}

在控制台命令行上输入命令：
\begin{enumerate}
  \item 单击控制台任意处以激活命令行。
  \item 键入命令。
  \item 单击 \texttt{de\_shell>} 提示按钮或按 Return。
\end{enumerate}

DC Explorer 会在 \cmd{de\_shell} 与控制台日志视图中回显命令输出，包括处理消息以及警告或错误消息。

要显示某 \cmd{de\_shell} 命令可用的选项，在命令行上输入命令名与 \opt{-help}。例如，查看 \cmd{compile\_exploration} 的选项：
\begin{lstlisting}
de_shell> compile_exploration -help
\end{lstlisting}

输入命令、选项或文件名时，可在键入足够多字符以唯一确定名称后按 Tab 键，由 DC Explorer 补全其余字符。
若已键入字符可对应多个名称，工具会列出候选名称，可用方向键与 Enter 键选择。

要从命令行界面的输出中复用命令，可选中相应文本，将指针移到 \cmd{de\_shell} 命令行，再单击鼠标中键粘贴。

\subsection{中断或终止命令处理}
\subsubsection*{Interrupting or Terminating Command Processing}

若为命令输入了错误选项或输入了错误命令，可中断命令处理并仍留在 \cmd{de\_shell} 中。
要中断或终止命令，按 Ctrl-C。

某些命令与进程（例如 \cmd{update\_timing}）无法中断。要停止这些命令或进程，必须在系统级终止 \cmd{de\_shell}。
终止进程或 shell 时不会保存任何数据。

按 Ctrl-C 时请注意：
\begin{itemize}
  \item 若正在处理脚本文件且您中断了其中某条命令，则脚本处理被中断，后续脚本命令不再执行。
  \item 若在命令响应中断之前按了三次 Ctrl-C，\cmd{de\_shell} 会被中断并以如下消息退出：
\begin{lstlisting}
Information: Process terminated by interrupt.
\end{lstlisting}
        少数例外行为记录在相应命令的 man 页中。
\end{itemize}

\subsection{在 de\_shell 中打开或关闭 GUI}
\subsubsection*{Opening or Closing the GUI in de\_shell}

启动命令行界面时，Linux shell 中会出现 \texttt{de\_shell>} 提示符。启动命令行界面后，可使用 Design Vision GUI。

Design Vision 在 Linux 的 X Window 环境下运行。打开 GUI 前，请确保 \texttt{\$DISPLAY} 环境变量设置为您 Linux 系统显示的名称。

若以无 GUI 方式启动 \cmd{de\_shell}，可在 \texttt{de\_shell>} 提示符下输入 \cmd{gui\_start} 打开 GUI。例如：
\begin{lstlisting}
de_shell> gui_start
\end{lstlisting}
也可在启动 DC Explorer 时对 \cmd{de\_shell} 指定 \opt{-gui} 选项以打开 GUI。
要从 \cmd{de\_shell} 会话使用 GUI，必须持有 Design Vision 许可证。

在会话中的任意时刻，都可打开或关闭 GUI 而不退出 DC Explorer。
例如，若需节省系统资源，可关闭 GUI 而让 DC Explorer 以命令行界面继续运行。

要不退出 Design Vision 而关闭 GUI，可任选其一：
\begin{itemize}
  \item 选择 File $>$ Close GUI
  \item \texttt{de\_shell> gui\_stop}
\end{itemize}

打开 GUI 时，它会读取一组名为 \texttt{.synopsys\_dv\_gui.tcl} 的 setup 文件。
可用这些文件执行 GUI 专用设置。\texttt{.synopsys\_dv\_gui.tcl} 中的设置会覆盖 \texttt{.synopsys\_dc.setup} 中的设置。
有关 \texttt{.synopsys\_dv\_gui.tcl} 的更多信息，请参阅 \textit{Design Vision User Guide}。

除读取 setup 文件外，GUI 还会从主目录中的 \texttt{\string~/.synopsys\_dc\_gui/preferences.tcl} 加载首选项与视图设置。
您不应编辑该文件。默认系统首选项面向最优工具运行，对多数设计适用。
但必要时可在会话中通过 Application Preferences 对话框更改 GUI 首选项。

可设置控制 GUI 窗口文本外观的首选项，以及选择或交互操作相关命令是否出现在会话日志中。
还可设置各类全局、原理图视图与版图视图的默认控件。

更多信息见 \textit{Design Vision User Guide} 中的 ``Setting GUI Preferences'' 主题。

\subsection{在命令行获取帮助}
\subsubsection*{Getting Help on the Command Line}

使用 DC Explorer 时可利用下列联机信息资源：
\begin{itemize}
  \item 命令帮助：指定 \cmd{de\_shell} 命令所用选项与参数的列表，显示在 shell 中；GUI 打开时也显示在控制台日志视图中
  \item Man 页：显示在 shell 中；GUI 打开时也显示在控制台日志视图中
  \item Man 页查看器：使用 GUI 时显示命令、变量与错误消息的 man 页
\end{itemize}

要显示 \cmd{de\_shell} 命令的简要说明：
\begin{lstlisting}
de_shell> help command_name
\end{lstlisting}

要显示命令选项与参数列表：
\begin{lstlisting}
de_shell> command_name -help
\end{lstlisting}

要在 \cmd{de\_shell}（以及 GUI 打开时的控制台日志视图）中显示 man 页：
\begin{lstlisting}
de_shell> man command_or_variable_name
\end{lstlisting}

要在 GUI man 页查看器中显示 man 页，可任选其一：
\begin{itemize}
  \item 在 GUI 控制台命令行上输入 \cmd{man} 后跟命令或变量名：
\begin{lstlisting}
de_shell> man command_or_variable_name
\end{lstlisting}
  \item 在 \cmd{de\_shell} 或 GUI 控制台命令行上输入 \cmd{gui\_show\_man\_page} 后跟命令或变量名：
\begin{lstlisting}
de_shell> gui_show_man_page command_or_variable_name
\end{lstlisting}
\end{itemize}

要在 GUI man 页查看器中交互浏览 man 页：
\begin{enumerate}
  \item 选择 Help $>$ Man Pages。\\
        Man 页查看器打开。
  \item 选择要查看的 man 页类型：Commands、Variables 或 Messages。\\
        随即出现 man 页列表。
  \item 选择要查看的 man 页。
\end{enumerate}

\subsection{Setup 文件}
\subsubsection*{Setup Files}

调用 DC Explorer 时，它会自动执行三个 setup 文件中的命令。
这些文件同名，均为 \texttt{.synopsys\_dc.setup}，但位于不同目录。
文件中包含初始化参数与变量、声明设计库等命令。

DC Explorer 按下列顺序从三个目录读取三个 \texttt{.synopsys\_dc.setup} 文件：
\begin{enumerate}
  \item Synopsys 根目录
  \item 您的主目录（home directory）
  \item 当前工作目录（即您调用 DC Explorer 的目录）
\end{enumerate}

表~\ref{tab:dce-setup-files} 说明这三个 setup 文件的作用。

\begin{table}[htbp]
\centering
\caption{Setup 文件}
\label{tab:dce-setup-files}
\begin{tabularx}{\textwidth}{@{}l l X@{}}
\toprule
\textbf{文件} & \textbf{位置} & \textbf{作用} \\
\midrule
系统级 \texttt{.synopsys\_dc.setup} &
Synopsys 根目录（\texttt{\$SYNOPSYS/admin/setup}） &
包含由 Synopsys 定义的系统变量以及面向站点全部用户的一般 DC Explorer 设置。仅系统管理员可修改此文件。注意：\texttt{\$SYNOPSYS} 为 DC Explorer 安装目录路径。 \\
用户定义 \texttt{.synopsys\_dc.setup} &
您的主目录 &
包含定义 DC Explorer 工作环境偏好的变量。此文件中的变量覆盖系统级 setup 文件中的对应变量。 \\
设计专用 \texttt{.synopsys\_dc.setup} &
启动 DC Explorer 的工作目录 &
包含影响该目录中全部设计优化的项目或设计专用变量。要使用此文件，必须从此目录调用 DC Explorer。此文件中定义的变量覆盖用户定义与系统级 setup 文件中的对应变量。 \\
\bottomrule
\end{tabularx}
\end{table}

示例~\ref{ex:synopsys-dc-setup} 给出 \texttt{.synopsys\_dc.setup} 文件示例。

\begin{lstlisting}[caption={.synopsys\_dc.setup 文件},label={ex:synopsys-dc-setup}]
# Define the target library, symbol library,
# and link libraries
set_app_var target_library lsi_10k.db
set_app_var synthetic_library dw_foundation.sldb
set_app_var link_library "* $target_library $synthetic_library"
set_app_var search_path [concat $search_path ./src]
set_app_var designer "Your Name"
# Define aliases
alias h history
alias rc "report_constraint -all_violators"
\end{lstlisting}

\subsection{在日志文件中查找会话信息}
\subsubsection*{Finding Session Information in the Log Files}

可在下列日志文件中查找会话信息（例如已处理的 \cmd{de\_shell} 命令与已访问的文件）：
\begin{itemize}
  \item 命令日志文件（Command Log Files）
  \item 编译日志文件（Compile Log Files）
  \item 文件名日志文件（Filename Log Files）
\end{itemize}

\subsubsection{命令日志文件}
\subsubsection*{Command Log Files}

命令日志文件记录 DC Explorer 处理的 \cmd{de\_shell} 命令，包括 setup 文件命令与变量赋值。
默认情况下，DC Explorer 将命令日志写入调用 \cmd{de\_shell} 的目录中名为 \texttt{command.log} 的文件。

可在 \texttt{.synopsys\_dc.setup} 文件中用 \varname{sh\_command\_log\_file} 变量更改 \texttt{command.log} 的文件名。
应在启动 DC Explorer 之前更改这些变量。
若用户定义或项目专用的 \texttt{.synopsys\_dc.setup} 不含该变量，DC Explorer 会自动创建 \texttt{command.log} 文件。

每次 DC Explorer 会话都会覆盖命令日志文件。要保存命令日志，请移动或重命名它。可用命令日志文件来：
\begin{itemize}
  \item 为特定综合策略生成脚本
  \item 记录设计探索过程
  \item 记录所遇到的问题
\end{itemize}

\subsubsection{编译日志文件}
\subsubsection*{Compile Log Files}

每次编译设计时，DC Explorer 会生成 ASCII 日志、HTML 日志文件以及
\texttt{.DE\_log\_snapshot\_date\_process\_id} 目录：
\begin{itemize}
  \item \textbf{ASCII 日志}\\
        该日志在屏幕上显示每次 DC Explorer 运行的输出（例如已处理命令与错误消息），便于快速查看与调试。
        日志末尾的摘要显示有多少消息被重定向到 \texttt{.DE\_log\_snapshot\_date\_process\_id} 目录下各具体日志文件。

  \item \textbf{HTML 日志文件}\\
        该文件位于当前工作目录下，以 HTML 格式包含 ASCII 日志的完整内容以及全部被重定向的消息。
        文件末尾的摘要表按消息 ID 分组给出所有消息出现情况的概览；单击消息 ID 链接即可显示该消息。
        默认文件名为 \texttt{default.html}。要更改文件名，请设置 \varname{de\_log\_html\_filename} 变量。例如：
\begin{lstlisting}
de_shell> set_app_var de_log_html_filename test1.log.html
\end{lstlisting}
\end{itemize}

\begin{noteBox}
必须安装 Python 编程语言才能生成 HTML 日志文件。下载信息见：\url{https://www.python.org/download}

默认情况下，工具将非关键警告与信息消息从 ASCII 日志重定向到 HTML 日志文件。
要禁用重定向，将 \varname{de\_log\_redirect\_enable} 变量设为 \texttt{false}。
\end{noteBox}

\begin{itemize}
  \item \textbf{\texttt{.DE\_log\_snapshot\_date\_process\_id} 目录}\\
        该目录存储从 ASCII 日志重定向出的信息与警告消息。错误消息显示在 ASCII 日志中，并保存在 HTML 日志文件中。
        该目录包含下列日志文件：
        \begin{itemize}
          \item \texttt{read\_design\_log}
          \item \texttt{read\_library\_log}
          \item \texttt{timing\_message\_log}
          \item \texttt{power\_message\_log}
          \item \texttt{ungroup\_hierarchies\_log}
          \item \texttt{register\_removal\_log}
          \item \texttt{optimization\_message\_log}
        \end{itemize}
\end{itemize}

\subsubsection{文件名日志文件}
\subsubsection*{Filename Log Files}

DC Explorer 将已读取文件的名称写入调用 \cmd{de\_shell} 的目录中的文件名日志文件。
若 DC Explorer 异常终止，可用该文件识别复现错误所需的数据文件。
要指定文件名日志文件的名称，在 \texttt{.synopsys\_dc.setup} 中设置 \varname{filename\_log\_file} 变量。

\subsection{运行 Tcl 脚本}
\subsubsection*{Running Tcl Scripts}

可用 Tcl 脚本来完成例行、重复或复杂任务。将一系列 \cmd{de\_shell} 命令放入文本文件即可创建命令脚本文件。
脚本文件中可执行任意 \cmd{de\_shell} 命令。要从 \cmd{de\_shell} 命令行运行脚本文件，输入：
\begin{lstlisting}
de_shell> source file_name
\end{lstlisting}

在 Tcl 中，行首的井号（\texttt{\#}）表示注释。例如：
\begin{lstlisting}
# This is a comment
\end{lstlisting}

要在 GUI 中运行脚本，使用 File $>$ Execute Scripts 对话框。

更多信息另见 \textit{Using Tcl With Synopsys Tools} 文档。

\subsection{退出 DC Explorer}
\subsubsection*{Exiting DC Explorer}

可随时退出 DC Explorer 并返回操作系统。默认情况下，\cmd{de\_shell} 将会话信息保存在 \texttt{command.log} 文件中。
但若在启动 DC Explorer 之后更改 \varname{sh\_command\_log\_file} 的文件名，会话信息可能会丢失。
此外，\cmd{de\_shell} 不会自动保存已加载到内存中的设计。若要在退出前保存这些设计，请使用 \cmd{write\_file} 命令。例如：
\begin{lstlisting}
de_shell> write_file -format ddc -hierarchy -output my_design.ddc
\end{lstlisting}

要退出 \cmd{de\_shell}，可任选其一：
\begin{itemize}
  \item 输入 \cmd{quit}。
  \item 输入 \cmd{exit}。
  \item 若以交互模式运行且工具正忙，按 Ctrl-d。
\end{itemize}

要在 Design Vision GUI 中退出 DC Explorer，选择 File $>$ Exit。

% ============================================================
\section{许可证}
\subsection*{Licensing}

下列各设计流程适用特定许可证要求：
\begin{itemize}
  \item DC Explorer
  \begin{itemize}
    \item RTL-Exploration 许可证
    \item DC-Explorer-Shell 许可证
  \end{itemize}
  \item DC Explorer floorplan 探索
  \begin{itemize}
    \item 一套 DC Explorer 许可证、一份 DC-Extension 许可证以及一份 ICC-DP 许可证；\\
          或者
    \item 两套 DC Explorer 许可证与两份 DC-Extension 许可证
  \end{itemize}
  \item Pre-DFT 设计规则检查
  \begin{itemize}
    \item TestMAX DFT 许可证
  \end{itemize}
\end{itemize}

Design Vision（GUI）与 DC-Extension 许可证会按需自动检出。

\begin{noteBox}
使用多核处理时，每两个核需要一份许可证。要移除多核处理期间额外检出的许可证，使用 \cmd{remove\_license} 命令。
\end{noteBox}

本节包括以下主题：
\begin{itemize}
  \item 列出正在使用的许可证
  \item 获取许可证
  \item 释放许可证
  \item 启用许可证排队
\end{itemize}

\subsection{列出正在使用的许可证}
\subsubsection*{Listing the Licenses in Use}

要查看当前已检出的许可证，使用 \cmd{list\_licenses} 命令。例如：
\begin{lstlisting}
de_shell> list_licenses
Licenses in use:
          DC-Explorer-Shell
          Design-Vision
          RTL-Exploration
1
\end{lstlisting}

要显示哪些许可证已被检出，使用 \cmd{license\_users} 命令。例如：
\begin{lstlisting}
de_shell> license_users
jack@eng1 RTL-Exploration
john@eng2 RTL-Exploration, Design-Vision
2 users listed.
\end{lstlisting}

\subsection{获取许可证}
\subsubsection*{Obtaining Licenses}

调用 DC Explorer 时，Synopsys Common Licensing 软件会自动检出相应许可证。
例如，若读入 HDL 设计描述，Synopsys Common Licensing 会为相应的 HDL Compiler 检出许可证。

可在 \cmd{de\_shell} 或 GUI 中获取许可证。要从 \cmd{de\_shell} 会话启动 GUI，必须持有 Design Vision 许可证。

\begin{itemize}
  \item \textbf{使用 \cmd{de\_shell}}\\
        若已知所需工具与接口，可用 \cmd{get\_license} 命令检出这些许可证，以确保准备使用时各许可证可用。
        默认情况下，每个功能仅检出一份许可证。许可证检出后一直保持，直到您释放它或退出 \cmd{de\_shell}。

        例如，下列命令为 HDL Compiler 功能检出一份许可证：
\begin{lstlisting}
de_shell> get_license HDL-Compiler
\end{lstlisting}
        对于可能需要多份许可证的多核处理，可对 \cmd{get\_license} 使用 \opt{-quantity} 选项指定所需许可证总数。
        若已检出部分许可证，该命令仅再获取使总数达到指定数量所需的额外许可证。

  \item \textbf{使用 GUI}
  \begin{itemize}
    \item 要显示当前许可证信息，选择 File $>$ Licenses。\\
          出现 Application Licenses 对话框。Allocated Licenses 列表显示当前正在使用的许可证；Available Licenses 列表显示其他可用许可证。
    \item 要检出额外许可证，选择 File $>$ Licenses，在 Available Licenses 列表中单击许可证名，再单击 Allocate。\\
          若有可用许可证，工具会检出一份副本；若许可证均已被占用，则显示错误消息。
  \end{itemize}
\end{itemize}

\begin{seeAlsoBox}
\begin{itemize}
  \item 列出正在使用的许可证
  \item 释放许可证
\end{itemize}
\end{seeAlsoBox}

\subsection{释放许可证}
\subsubsection*{Releasing Licenses}

可在 \cmd{de\_shell} 或 GUI 中释放许可证。要从 \cmd{de\_shell} 会话启动 GUI，必须持有 Design Vision 许可证。

要释放已检出到您名下的许可证，使用 \cmd{remove\_license} 命令。例如：
\begin{lstlisting}
de_shell> remove_license HDL-Compiler
\end{lstlisting}
可用 \opt{-keep} 选项指定命令完成后为每个功能保留的许可证数量。
若不指定该选项，运行 \cmd{remove\_license} 时会释放每个功能的全部许可证。

要在 GUI 中释放许可证：
\begin{enumerate}
  \item 选择 File $>$ Licenses。
  \item 在 Allocated Licenses 列表中选择一份许可证。
  \item 单击 Release。
\end{enumerate}

\begin{seeAlsoBox}
\begin{itemize}
  \item 列出正在使用的许可证
  \item 获取许可证
\end{itemize}
\end{seeAlsoBox}

\subsection{启用许可证排队}
\subsubsection*{Enabling License Queuing}

DC Explorer 具有许可证排队功能：当全部许可证均在使用时，允许应用程序等待许可证变为可用。
要启用该功能，将环境变量 \varname{SNPSLMD\_QUEUE} 设为 \texttt{true}。将显示如下消息：
\begin{lstlisting}
Information: License queuing is enabled. (DCSH-18)
\end{lstlisting}

启用许可证排队后，可能出现您持有许可证 L1 却在等待 L2，而另一用户持有 L2 却在等待 L1 的情况；此时双方将一直等待。

为防止此类情况，请使用 \varname{SNPS\_MAX\_WAITTIME} 与 \varname{SNPS\_MAX\_QUEUETIME} 环境变量。
在使用这两个变量之前，必须先将 \varname{SNPSLMD\_QUEUE} 设为 \texttt{true}。
\begin{itemize}
  \item \varname{SNPS\_MAX\_WAITTIME} 指定您所需第一份关键许可证的最大等待时间。
  \item \varname{SNPS\_MAX\_QUEUETIME} 指定在同一 \cmd{de\_shell} 进程中检出后续许可证的最大等待时间。
        在成功检出第一份许可证以启动 \cmd{de\_shell} 之后使用该变量。
\end{itemize}

当您将设计跑过综合流程时，排队功能可能显示其他状态消息，例如：
\begin{lstlisting}
Information: Started queuing for feature 'HDL-Compiler'. (DCSH-15)
Information: Still waiting for feature 'HDL-Compiler'. (DCSH-16)
Information: Successfully checked out feature 'HDL-Compiler'. (DCSH-14)
\end{lstlisting}

% ============================================================
\section{使用多核技术}
\subsection*{Using Multicore Technology}

DC Explorer 中的多核技术允许使用多个核以改善工具运行时间。
在综合期间，该技术可将大型优化任务划分为较小任务，以便在多个核上处理。

\subsection{启用多核处理}
\subsubsection*{Enabling Multicore Processing}

要在 DC Explorer 中启用多核处理，使用 \cmd{set\_host\_options} 命令。
例如，要使工具使用六个核运行进程，输入：
\begin{lstlisting}
de_shell> set_host_options -max_cores 6
\end{lstlisting}
所有 \cmd{compile\_exploration} 命令选项均支持使用多核进行优化。

启用多核处理后，日志文件会包含类似如下的信息消息：
\begin{lstlisting}
Information: Running optimization using a maximum of 6 cores. (OPT-1500)
\end{lstlisting}
为获得最大效率，\opt{-max\_cores} 设置应不大于实际可用核数。

设置 \varname{compile\_adjust\_max\_processes\_used} 变量时，工具按该变量的下列取值运行多核进程：
\begin{itemize}
  \item \texttt{1}：默认情况下，若 \texttt{overcommit\_memory} 参数设为 2，工具在多进程期间调整核数。
        若资源不足以并发运行新进程：
        \begin{itemize}
          \item \cmd{compile} 或 \cmd{compile\_ultra} 会调整核数以继续，而不是终止进程。
                但仅当至少一个新进程成功时，命令才会继续执行。
          \item \cmd{parallel\_execute} 会减少要用的进程数；若第一个新进程失败，则在主进程中继续执行。
          \item \cmd{redirect -bg} 即使新进程失败，也会在主进程中继续执行。
        \end{itemize}
  \item \texttt{2}：无论 \texttt{overcommit\_memory} 参数设置如何，都调整核数。
  \item \texttt{0}：不调整核数。
\end{itemize}

\begin{seeAlsoBox}
调整多核处理期间的核数（下文）。
\end{seeAlsoBox}

\subsection{报告运行时间}
\subsubsection*{Reporting Runtime}

要报告总体编译挂钟时间（wall clock time），运行 \cmd{report\_qor} 命令，如下例所示。
该命令报告运行 \cmd{compile\_exploration} 的挂钟时间。
\begin{lstlisting}
de_shell > report_qor
****************************************
Report : qor
Design : TEST_TOP
...
****************************************
...
  Hostname: machine
  Compile CPU Statistics
  -----------------------------------------
  ...
  -----------------------------------------
  Overall Compile Time:              631.32
  Overall Compile Wall Clock Time:   288.11
\end{lstlisting}

所报告的挂钟时间与下列脚本报告的时间类似：
\begin{lstlisting}
de_shell> set_host_options -max_cores 2
de_shell> set pre_compile_clock [clock seconds]
de_shell> compile_exploration
de_shell> set post_compile_clock [clock seconds]
de_shell> set diff_clock [expr $post_compile_clock - $pre_compile_clock]
\end{lstlisting}

用多核优化衡量运行时间加速时，应使用进程的挂钟时间。CPU 时间不是多核运行时间加速的正确指标。

\begin{seeAlsoBox}
报告结果质量（QoR）。
\end{seeAlsoBox}

\subsection{调整多核处理期间的核数}
\subsubsection*{Adjusting the Number of Cores During Multicore Processing}

要在 \texttt{overcommit\_memory} 参数设为 2 时，默认允许 Design Compiler 工具在多核处理期间调整核数，请使用 \varname{compile\_adjust\_max\_processes\_used} 变量。
该变量允许工具通过发出下列消息继续执行：
\begin{lstlisting}
UIO-306 (Warning) There is insufficient memory to launch all jobs,
 reducing the number of
parallel processes to %d for this session. Tool performance may be
 impacted. Please set
overcommit_memory to 0.
\end{lstlisting}

若未使用 \varname{compile\_adjust\_max\_processes\_used} 变量且将 \texttt{overcommit\_memory} 参数设为 2，操作系统可能对多进程程序错误地使用 copy-on-write 优化策略，为所有进程增加虚拟内存，即使有空闲内存也不分配。
这会导致正在运行的进程或创建新进程时被阻塞，并使工具崩溃，发出如下消息：
\begin{lstlisting}
OPT-1603 (error) Insufficient virtual memory. Please increase swap space,
or reduce other processes on this host.
\end{lstlisting}

\begin{seeAlsoBox}
启用多核处理（上文）。
\end{seeAlsoBox}

% ============================================================
\section{并行运行命令}
\subsection*{Running Commands in Parallel}

在脚本中串行执行检查或报告命令可能占用总体运行时间的很大一部分。
为改善运行时间并生成相同报告，可使用并行命令执行来并行运行这些命令。
\begin{itemize}
  \item \cmd{parallel\_execute} 命令
  \item 许可证要求
  \item 多核处理
  \item 支持的命令
\end{itemize}

\subsection{parallel\_execute 命令}
\subsubsection*{The parallel\_execute Command}

要启用并行命令执行，在脚本中用 \cmd{parallel\_execute} 命令列出要并行执行的检查或报告命令。
若列表中的某命令需要此前未执行过的时序更新，该命令会在并行命令执行期间自动调用 \cmd{update\_timing}。

例如，下列脚本先更新时序信息，然后使用八个核并行执行 \cmd{report\_timing}、\cmd{report\_qor}、\cmd{report\_cell} 与 \cmd{report\_area}：
\begin{lstlisting}
set_host_options -max_cores 8
update_timing
parallel_execute [list \
  "report_timing > $mylogfile" \
  "report_qor >> $mylogfile" \
  "report_cell" \
  "report_area"]
\end{lstlisting}

下列脚本在命令字符串中对选项参数与日志文件使用变量（包括 \texttt{\$MAX}、\texttt{\$NWORST}、\texttt{\$cstr\_log} 与 \texttt{\$qor\_log}）：
\begin{lstlisting}
set MAX 10000; set_app_var NWORST 10
set rpt_const_options "-all_violators"
set cstr_log "rpt_cstr.log"
set qor_log "rpt_qor.log"
parallel_execute [list \
   "report_timing -max $MAX -nworst $NWORST > rpt_tim.log" \
   "report_constraints $rpt_const_options > $cstr_log" \
   "report_qor > $qor_log"]
\end{lstlisting}

\subsection{许可证要求}
\subsubsection*{License Requirements}

对使用多核处理的并行命令执行，适用与串行运行命令相同的许可证要求。
默认情况下，每八个核需要一套 DC Explorer 许可证。

\subsection{多核处理}
\subsubsection*{Multicore Processing}

并行命令执行最多可使用八个核。一次并行执行的最大命令数由 \cmd{set\_host\_options} 的 \opt{-max\_cores} 选项决定。
若指定大于八的数字，工具发出警告并用可用核数覆盖该设置。
工具会阻塞 \cmd{de\_shell}，直到并行执行列表中运行时间最长的命令完成。

下列示例使用三个核运行并行命令执行：
\begin{lstlisting}
de_shell> parallel_execute [list \
   "report_cell" "report_timing" "report_area"]
Information: Running parallel report using a maximum of 3 cores.
(RPT-100)
\end{lstlisting}

下列示例为并行命令执行指定八个核：
\begin{lstlisting}
de_shell> set_host_options -max_cores 8
\end{lstlisting}

\subsection{支持的命令}
\subsubsection*{Supported Commands}

仅对报告与检查命令使用并行命令执行。要查明支持哪些命令，对 \cmd{parallel\_execute} 指定 \opt{-list\_all} 选项。
若指定不支持的命令，工具会跳过该命令并发出类似如下的警告：
\begin{lstlisting}
de_shell> parallel_execute [list report_libcell_subset]
Warning: 'report_libcell_subset' report command can't run in
 parallel_execute mode.
(RPT-106)
\end{lstlisting}

\subsection{在设计流程中}
\subsubsection*{In the Design Flow}

下图说明如何在 DC Explorer 设计流程中并行运行检查与报告命令以改善运行时间。

\figplaceholder{Figure 2: Parallel Command Execution in the DC Explorer Design Flow}{DC Explorer 设计流程中的并行命令执行}{fig:parallel-exec}

\begin{seeAlsoBox}
并行运行命令（本节上文）。
\end{seeAlsoBox}

% ============================================================
\section{在后台运行命令以改善运行时间}
\subsection*{Running Commands in the Background to Improve Runtime}

要在 DC Explorer 中改善运行时间，可对 \cmd{redirect} 命令使用 \opt{-bg} 选项，在前台运行其他命令的同时在后台运行命令。
可运行单条命令、\cmd{source} 一个 Tcl 脚本以并行运行若干命令，或运行一个 Tcl 过程。

要用 \cmd{redirect -bg} 指定是在后台运行只读命令还是非只读命令，对 \cmd{enable\_redirect\_bg\_commands} 使用下列选项：
\begin{itemize}
  \item \opt{-all}（默认）：允许工具用 \cmd{redirect -bg} 在后台运行只读与非只读命令。
  \item \opt{-read\_only}：仅用 \cmd{redirect -bg} 在后台运行只读命令；所有其他（非只读）命令在前台运行。
        若对 \opt{-read\_only} 指定了非只读命令，工具发出 BGE-011 信息消息。
\end{itemize}

运行后台作业时必须指定要用的核数。\cmd{redirect -bg} 可用的核数从 \cmd{set\_host\_options} 设置的核数中扣除。
使用 \cmd{redirect} 运行后台作业时，\cmd{de\_shell} 立即返回以执行下一条命令。

下列示例将 Tcl 脚本重定向到后台运行，并使用 GUI 进行分析：
\begin{lstlisting}
set_host_options -max_cores 8
read_verilog
read_sdc
...
compile_exploration
redirect -bg -max_cores 3 -file y.out {source y.tcl}
remote_execute_source1.log {source bg_script.tcl}
   Information: redirect -bg with max_cores 3 started. The maximum number
   of cores available in parent is reduced to 5. (BGE-004)
start_gui
\end{lstlisting}

同一时间最多只能有两个活动的后台作业。若再重定向第三条命令或脚本，工具会等待先前命令完成后再启动新作业。

支持的命令仅包括报告、分析与设计导出命令。会更改网表或其他数据的命令不受支持。

也可在 \cmd{redirect -bg} 命令内部运行 \cmd{parallel\_execute}，以并行运行一组检查与报告命令。

\subsection{报告用 redirect 提交的作业}
\subsubsection*{Reporting Jobs Submitted With the redirect Command}

要报告您已用 \cmd{redirect -bg} 提交到后台运行的全部作业，使用 \cmd{report\_background\_jobs} 命令。
该命令报告已完成作业以及当前正在后台运行的作业，如下例所示：
\begin{lstlisting}
de_shell> report_background_jobs

JOB 'redirect -bg  -file {background.log} source bg_script.tcl
 -max_cores 4 ' completed
JOB(pid:13010) 'redirect -bg  -file {background_1.log} source
 bg_script_1.tcl  -max_cores 3 ' is running
\end{lstlisting}

要省略已完成作业，对 \cmd{report\_background\_jobs} 使用 \opt{-reset} 选项。
